\section{Compare a second correct instruction strategy}
The alternative bit implementation converts the retained low bit into a zero-or-all-ones
mask, ANDs it with the polynomial, and XORs the result after the shift. It performs the
feedback XOR even when that masked value is zero. This removes the per-bit conditional
branch from the illustrated update. The alternative table path uses a signed word load
followed by explicit zero extension instead of an unsigned word load.

\begin{table}[htbp]
  \centering
  \caption{Retired instructions for two independently checked strategies.}
  \tagpdfsetup{table/header-rows={1,2}}
  \begin{tabular}{lrrrr}
    \toprule
                   & \multicolumn{2}{c}{Bit update} & \multicolumn{2}{c}{Prepared table}                                    \\
    Dataset        & Branch                         & Mask                               & Unsigned load & Load then extend \\
    \midrule
    Empty          & 12                             & 12                                 & 5             & 5                \\
    One byte       & 84                             & 85                                 & 19            & 20               \\
    Check message  & 631                            & 669                                & 131           & 140              \\
    Zero bytes     & 4415                           & 4684                               & 901           & 965              \\
    All octets     & 17709                          & 18700                              & 3589          & 3845             \\
    Strided octets & 282498                         & 299020                             & 57349         & 61445            \\
    \bottomrule
  \end{tabular}
\end{table}

The alternative initializer retires 18444 instructions. Both implementations produce
the same table words, raw states and finalized checksums under the tested contract.
The additional zero-extension instruction accounts for one extra instruction per byte
in the alternative table path. The mask bit loop executes a fixed feedback sequence,
while the branch path skips the XOR when feedback is zero.

These results measure simulator instruction work. They establish neither elapsed-time
ratios nor hardware cycle counts. A branch, a table load and a syscall need not have
the same cost in a host implementation or a physical processor. File open/read/close
latency is outside the kernel measurements. If elapsed time matters, measure the complete
file operation separately and record input size, buffer size, environment and repetitions.

\subsection*{Exercise: Charge setup to the work that uses it}
For the six datasets, add one table-construction cost to the prepared-table update.
Then consider one table reused across ten messages. State which setup boundary you
chose and which conclusions depend on data, rather than reporting an unconditional
speedup for every file.
